\section{Mixed Gamma derivatives and exact harmonic tail identities}
\label{mx:section}
The harmonic coefficient calculus has a useful companion at a double
zero of a hypergeometric residue. The mixed spectral continuation
\cite{mx:report} identifies a normally convergent two-parameter Gamma
identity and a weighted square of trigamma differences. We prove the
needed first resonance directly from the classical Rogers--Dougall sum
\cite{mx:DLMFDougall}, rather than assuming the entire pending
all-resonance programme. Its mixed derivatives produce ordinary harmonic
identities. Differentiating the resulting square identity also suggests
an all-orders zeta-tail family; an independent Euler partial-fraction
proof of that family closes the loop between experiment and exact algebra.

\subsection{A normally convergent double-zero germ}
Put $\eta=a/2$, $x_n=n+\eta$, and $A=a-1$. Define
\begin{align}
 Q_n(a;p,q)&=
 \frac{\Gamma(n+a)\Gamma(n+1+p)\Gamma(n+1+q)\Gamma(n+a-p-q)}
 {\Gamma(n+1)\Gamma(n+a-p)\Gamma(n+a-q)\Gamma(n+1+p+q)},
 \label{mx:Q}\\
 B_a(p,q)&=\psi(a/2)+\tfrac12\{
 \psi(2)-\psi(1+p)-\psi(1+q)-\psi(a-p-q)\}.
 \label{mx:B}
\end{align}
These are holomorphic for $\Re a>0$ and $p,q$ sufficiently small
on each compact $a$-set. Both $Q_n(a;p,0)$ and $Q_n(a;0,q)$ are one.

\begin{theorem}[The first mixed resonance]\label{mx:generator}
For every real $a>0$, and sufficiently small complex $p,q$,
\begin{equation}\label{mx:generator-eq}
 \sum_{n\ge0}\left\{
 x_n\bigl(Q_n(a;p,q)-1\bigr)+\frac{(A-p-q)pq}{x_n}\right\}
 =-(A-p-q)pq\,B_a(p,q).
\end{equation}
The bracketed series converges locally normally, including after
division by $pq$ at either axis. Every fixed parameter derivative
can be taken termwise. The identity extends holomorphically to the
corresponding germs over $\Re a>0$.
\end{theorem}
\begin{proof}
First take $a>1$, $b=1+p$, $c=1+q$ and $d_0=1+a-b-c>0$.
Introduce a separate spectral variable $u$ and the centered summand
\[
 r_n(u)=x_n\frac{\Gamma(n+a)\Gamma(n+b)\Gamma(n+c)\Gamma(n+d_0-u)}
 {\Gamma(n+1)\Gamma(n+1+a-b)\Gamma(n+1+a-c)\Gamma(n+b+c+u)}.
\]
For $0<\Re u<d_0$, the Rogers--Dougall formula with upper parameters
$a,1+a/2,b,c,d_0-u$ and lower parameters
$a/2,1+a-b,1+a-c,b+c+u$ gives
\begin{equation}\label{mx:dougall-normalized}
 \sum_{n\ge0}r_n(u)=
 G(u):=\frac{\Gamma(b)\Gamma(c)\Gamma(d_0-u)\Gamma(u)}
 {2\Gamma(d_0)\Gamma(b+u)\Gamma(c+u)}.
\end{equation}
The classical excess is $2u$; the factor $1/2$ follows from the
normalization $x_n=n+a/2$, not the usual Pochhammer summand.

Write
\[
 (\alpha_1,\alpha_2,\alpha_3,\alpha_4)
 =(\eta,b-\eta,c-\eta,d_0-u-\eta).
\]
The denominator shifts relative to $x_n$ are $1-\alpha_i$.
The shifted Stirling expansion \cite{hstruct:DLMF} and Bernoulli
reflection therefore give uniformly on compact parameter sets
\[
 r_n(u)=x_n^{-1-2u}\{1+C_1(u)x_n^{-2}+O(n^{-4})\},\qquad
 C_1(u)=-\tfrac13\sum_{i=1}^4 B_3(\alpha_i).
\]
Centering cancels the odd inverse powers. Thus
\begin{equation}\label{mx:subtracted}
 R_2(u)=\sum_{n\ge0}
 \{r_n(u)-x_n^{-1-2u}-C_1(u)x_n^{-3-2u}\}
 =G(u)-\zeta(1+2u,\eta)-C_1(u)\zeta(3+2u,\eta)
\end{equation}
extends holomorphically to $-2<\Re u<d_0$. The series bound is
$O(n^{-5-2\Re u})$ on compact subsets; Cauchy's formula gives
normal convergence for each fixed spectral or parameter derivative.
The equality first follows by summing the counterterms in the
ordinary convergence region, and then by continuation.

At $u=-1$, direct Bernoulli algebra gives
\begin{align}
 C_1(-1)&=\rho:=-(A-p-q)pq,\notag\\
 \tfrac12 C_1'(-1)+\zeta(-1,\eta)
 &=\tfrac12\{B_2(\eta-p-q)-B_2(\eta)\}
 =-\tfrac12(A-p-q)(p+q).
 \label{mx:bernoulli-cancellation}
\end{align}
For generic nonzero $p,q$, $G$ has residue $\rho/2$ and constant term
\[
 \frac\rho2\{\psi(2)-\psi(d_0+1)-\psi(p)-\psi(q)\}.
\]
The resonant Hurwitz term has constant coefficient
$C_1'(-1)/2-\rho\psi(\eta)$. Combining these constants with
\eqref{mx:bernoulli-cancellation}, and using
$\psi(1+p)=\psi(p)+1/p$ and its $q$ counterpart, yields
$R_2(-1)=\rho B_a(p,q)$. Since $r_n(-1)=x_nQ_n$, this is
\eqref{mx:generator-eq}. Holomorphy extends it to the spectral axes.

Finally all Gamma arguments in $Q_n$ and all digamma arguments in
$B_a$ stay pole-free for $\Re a>0$ and small $p,q$. The centered
bound remains $O(n^{-3})$ uniformly on every compact $a$-set.
Each Taylor coefficient is consequently holomorphic in that half-plane
and agrees with the asserted coefficient for $\Re a>1$; the identity
theorem extends it throughout. The bracket vanishes identically on
both axes. Cauchy's formula on a slightly larger polydisc transfers
its $O(n^{-3})$ bound to its holomorphic quotient by $pq$ and to
every fixed derivative, proving the divided-series assertion.
\end{proof}

\subsection{Finite mixed coefficient tables}
For $k\ge2$ set
$\Lambda_k(n;a)=(-1)^k\psi^{(k-1)}(n+a)-\psi^{(k-1)}(n+1)$.
Logarithmic differentiation of the eight Gamma factors gives
\begin{equation}\label{mx:logQ}
 \log Q_n=\sum_{r,s\ge1}\Lambda_{r+s}(n;a)\frac{p^rq^s}{r!s!}.
\end{equation}
Define the finite bivariate Bell coefficient
\[
 \mathcal B_{rs}(n;a)=r!s![p^rq^s]
 \exp\left(\sum_{i,j\ge1}\Lambda_{i+j}(n;a)\frac{p^iq^j}{i!j!}\right).
\]
Only $i\le r$, $j\le s$ and powers at most $\min(r,s)$ are needed.
For example
\[
 \mathcal B_{11}=\Lambda_2,\quad
 \mathcal B_{31}=\Lambda_4,\quad
 \mathcal B_{22}=\Lambda_4+2\Lambda_2^2,\quad
 \mathcal B_{33}=\Lambda_6+9\Lambda_2\Lambda_4+9\Lambda_3^2+6\Lambda_2^3.
\]
Write $B_{rs}=\partial_p^r\partial_q^sB_a(0,0)$. Then
\[
 B_{00}=\psi(a/2)+\tfrac12(1+\gamma-\psi(a)),
\]
and, for $r+s=k\ge1$,
\begin{equation}\label{mx:B-jets}
 B_{rs}=-\tfrac12\{
 \mathbf1_{s=0}\psi^{(r)}(1)+\mathbf1_{r=0}\psi^{(s)}(1)
                  +(-1)^k\psi^{(k)}(a)\}.
\end{equation}
Terms with a negative subscript are omitted below.
\begin{corollary}[All fixed mixed orders]\label{mx:mixed-table}
For $r,s\ge1$ the absolutely convergent single bracket
\begin{align}
 S_{rs}(a)=\sum_{n\ge0}\bigg\{x_n\mathcal B_{rs}(n;a)
 +\frac{A\mathbf1_{r=s=1}-2\mathbf1_{(r,s)=(2,1)}
                         -2\mathbf1_{(r,s)=(1,2)}}{x_n}\bigg\}
 \label{mx:mixed-sum}
\end{align}
equals
\begin{equation}\label{mx:mixed-value}
 -rs\{A B_{r-1,s-1}-(r-1)B_{r-2,s-1}-(s-1)B_{r-1,s-2}\}.
\end{equation}
When $r+s\ge4$ the counterterm is absent.
\end{corollary}
\begin{proof}
Differentiate the normally convergent identity
\eqref{mx:generator-eq}. Its counterterm is a polynomial of total
degree three, while Leibniz's rule on the right gives
\eqref{mx:mixed-value}. The quotient and normal-convergence
statements were proved before differentiation.
\end{proof}

\subsection{A trigamma square and its arithmetic specializations}
\begin{theorem}[Weighted trigamma square]\label{mx:trigamma-square}
For every real $a>0$,
\begin{equation}\label{mx:trigamma-square-value}
 \sum_{n\ge0}\left(n+\frac a2\right)
       \bigl(\psi'(n+a)-\psi'(n+1)\bigr)^2
 =\frac{a-1}{4}\psi''(a)+\frac{\psi'(a)-\zeta(2)}2
                         +\frac{3(a-1)}2\zeta(3).
\end{equation}
The ordinary series is locally normally convergent over $\Re a>0$,
where the square is algebraic.
\end{theorem}
\begin{proof}
The $(2,2)$ and $(3,1)$ rows of \eqref{mx:mixed-value} give
\begin{align*}
 \sum_nx_n(\Lambda_4+2\Lambda_2^2)
 &=2A\psi''(a)+4\psi'(a)-4\zeta(2),\\
 \sum_nx_n\Lambda_4
 &=\tfrac{3A}{2}\{\psi''(1)+\psi''(a)\}
                       +3\{\psi'(a)-\zeta(2)\}.
\end{align*}
Subtract and divide by two, using $\psi''(1)=-2\zeta(3)$.
All series in that calculation converge absolutely. Independently,
the polygamma expansion gives
$\psi'(n+a)-\psi'(n+1)=O(n^{-2})$ locally uniformly, so the
weighted square is $O(n^{-3})$, including at $a=1$.
\end{proof}
The half-integer recurrences reduce the cases $a=3/2$ and $a=1/2$ to
\begin{align}
 \sum_{n\ge0}(4n+3)
 \{2H_{2n+1}^{(2)}-H_n^{(2)}-\zeta(2)\}^2
 &=\zeta(2)-\zeta(3),\label{mx:half-square}\\
 \sum_{n\ge0}(4n+1)
 \{\zeta(2)-2H_{2n}^{(2)}+H_n^{(2)}\}^2
 &=\zeta(2)+\zeta(3).\label{mx:other-half-square}
\end{align}
Indeed $\psi'(3/2)=3\zeta(2)-4$,
$\psi''(3/2)=-14\zeta(3)+16$,
$\psi'(1/2)=3\zeta(2)$ and $\psi''(1/2)=-14\zeta(3)$.
The polygamma series express the respective differences as twice
the bracketed harmonic tails. At $a=2$,
\eqref{mx:trigamma-square-value} reduces to $\sum_{k\ge1}k^{-3}=\zeta(3)$.

\subsection{An all-orders zeta-tail identity with an independent proof}
Put $T_s(n)=\zeta(s)-H_n^{(s)}=\sum_{j>n}j^{-s}$ for $s\ge2$.
Differentiating \eqref{mx:trigamma-square-value} at $a=1$ suggests the
following finite combination at every order. We prove it independently
by Euler partial fractions and positive double sums, so its validity
does not depend on the hypergeometric continuation.
\begin{theorem}[All-order quadratic tail reduction]\label{mx:tail-family}
For every integer $k\ge2$,
\begin{align}
 &\sum_{n\ge0}\bigg[
 \left(n+\tfrac12\right)\sum_{j=1}^{k-1}(j+1)(k-j+1)
                 T_{j+2}(n)T_{k-j+2}(n)\notag\\
 &\hspace{30mm}-\tfrac12\sum_{j=1}^{k-2}(j+1)(k-j)
                 T_{j+2}(n)T_{k-j+1}(n)\bigg]
 =\frac{(k+1)(k+2)}4\zeta(k+2).
 \label{mx:tail-family-value}
\end{align}
Every series displayed on the left is ordinary and absolutely convergent.
\end{theorem}
\begin{proof}
For $a,b\ge3$, summing first over $0\le n<\min(l,m)$ gives
\begin{align}
 \sum_{n\ge0}(n+\tfrac12)T_a(n)T_b(n)
 &=\tfrac12\{\zeta(a+b-2)+\zeta(a,b-2)+\zeta(b,a-2)\},
 \label{mx:weighted-tail}\\
 \sum_{n\ge0}T_a(n)T_b(n)
 &=\zeta(a+b-1)+\zeta(a,b-1)+\zeta(b,a-1).
 \label{mx:ordinary-tail}
\end{align}
The finite sums of the weights are $\min(l,m)^2/2$ and $\min(l,m)$,
respectively. Positivity justifies Tonelli, and the resulting double
zeta sums converge. Here $\zeta(r,s)=\sum_{l>m\ge1}l^{-r}m^{-s}$.
Substituting these two formulas, using the symmetry in $j$ and $k-j$
or $k-1-j$, reduces \eqref{mx:tail-family-value} to
\begin{equation}\label{mx:Euler-weighted}
 \sum_{j=1}^{k-2}(j+1)\zeta(j+2,k-j)
          +2k\zeta(k+1,1)=\zeta(k+2).
\end{equation}
The diagonal single-zeta coefficient before this last reduction is
$(k^2+3k-2)/4$; the proposed value exceeds it by one.

For completeness the classical relation \eqref{mx:Euler-weighted}
has an elementary proof. Iterating
\[
 \frac1{x^r y^s}=
 \frac{x^{-(r-1)}y^{-s}+x^{-r}y^{-(s-1)}}{x+y}
\]
gives, for $x,y>0$,
\begin{equation}\label{mx:partial-fraction}
 \frac1{x^2y^k}
 =\sum_{i=0}^{1}\frac{\binom{k+i-1}{i}}{(x+y)^{k+i}x^{2-i}}
  +\sum_{i=0}^{k-1}\frac{i+1}{(x+y)^{2+i}y^{k-i}}.
\end{equation}
Summing over positive integer $x,y$ is justified by positivity. The
first sum becomes $\zeta(k,2)+k\zeta(k+1,1)$; the second becomes
$\sum_{i=0}^{k-1}(i+1)\zeta(i+2,k-i)$.
On the other hand, separating equal and unequal indices in the product
gives the stuffle identity
$\zeta(2)\zeta(k)=\zeta(2,k)+\zeta(k,2)+\zeta(k+2)$.
Subtracting the two expressions proves \eqref{mx:Euler-weighted},
including $k=2$ where the empty sum is zero. This proves the theorem.
\end{proof}
The first member is the simple ordinary identity
\begin{equation}\label{mx:tail-three-square}
 \boxed{\sum_{n\ge0}(2n+1)\{\zeta(3)-H_n^{(3)}\}^2
                         =\tfrac32\zeta(4).}
\end{equation}
The next member is
\[
 \sum_{n\ge0}\{12(n+\tfrac12)T_3(n)T_4(n)-2T_3(n)^2\}
                         =5\zeta(5).
\]

\begin{corollary}[Centered tail squares]\label{mx:centered-tails}
For every integer $s\ge3$,
\begin{equation}\label{mx:centered-tail-family}
 \sum_{k\ge1}k\{2T_s(k-1)-k^{-s}\}^2
 =2\zeta(2s-2)+4\zeta(s,s-2)-\zeta(2s-1).
\end{equation}
At $s=3$ this is the mixed-report evaluation
\begin{equation}\label{mx:centered-three}
 \boxed{\sum_{k\ge1}k
       \{2(\zeta(3)-H_{k-1}^{(3)})-k^{-3}\}^2
                   =3\zeta(4)-\zeta(5).}
\end{equation}
\end{corollary}
\begin{proof}
Expand the square and use \eqref{mx:weighted-tail},
\eqref{mx:ordinary-tail}, and
$\sum_{k\ge1}k^{1-s}T_s(k-1)=\zeta(s,s-1)+\zeta(2s-1)$.
The two $\zeta(s,s-1)$ terms cancel. For $s=3$,
\eqref{mx:Euler-weighted} at $k=2$ gives
$4\zeta(3,1)=\zeta(4)$, yielding \eqref{mx:centered-three}.
All rearranged sums converge absolutely.
\end{proof}
These identities show how a Gamma mixed-derivative experiment can expose
an exact Euler relation. The additional all-order tail and centered-tail
formulas are collective deductions from classical partial fractions;
global priority or independence of their period coordinates is not asserted.
